% FAC-EN-U0060
% Source: fac_body.tex lines 2195--2267.
% Authority: PDF physical pp. 59--61; printed pp. 254--256
% (Chapter III, no. 62 complete).
% French control T0035 restores four authority-visible membership signs.

\phantomsection\label{fac:III:3:no:62}
\textbf{62. Computation of $H^q_k(M)$ for certain modules $M$.}
Let $M$ be a graded $S$-module, and let $m\in M$ be a homogeneous
element of degree $0$. The system $(t_i^k.m)$ is a $0$-cocycle of
$C_k(M)$, which we shall denote by $\alpha^k(m)$ and identify with its
cohomology class. This gives a $K$-linear homomorphism
\phantomsection\label{fac:III:3:no:62:def:alpha-k}
\[
\alpha^k:M_0\longrightarrow H^0_k(M).
\]
Since $\alpha^h=\rho^h_k\circ\alpha^k$ when $h\geq k$, passage to the
limit gives a homomorphism
\phantomsection\label{fac:III:3:no:62:def:alpha}
\[
\alpha:M_0\longrightarrow H^0(M).
\]

We introduce two further notations.

\phantomsection\label{fac:III:3:no:62:def:generated-submodule}
If $P_0,\ldots,P_h$ are elements of $S$, we denote by
$(P_0,\ldots,P_h)M$ the submodule of $M$ consisting of the elements
\[
\sum_{i=0}^{i=h}P_i.m_i,\qquad m_i\in M.
\]
If the $P_i$ are homogeneous, this submodule is homogeneous.

\phantomsection\label{fac:III:3:no:62:def:colon-submodule}
If $P$ is an element of $S$ and $N$ a submodule of $M$, we denote by
$N:P$ the submodule of $M$ consisting of the elements $m\in M$ such
that $P.m\in N$. Evidently $N:P\supset N$; if $N$ and $P$ are
homogeneous, then $N:P$ is homogeneous.

With these notations fixed, we have:

\phantomsection\label{fac:III:3:no:62:prop:1}
\textsc{Proposition 1.} \textit{Let $M$ be a graded $S$-module and let
$k$ be an integer $\geq0$. Suppose that, for every $i\in I$,}
\phantomsection\label{fac:III:3:no:62:prop:1:hypothesis}
\[
(t_0^k,\ldots,t_{i-1}^k)M:t_i^k
=(t_0^k,\ldots,t_{i-1}^k)M.
\]
\textit{Then:}

\phantomsection\label{fac:III:3:no:62:prop:1:item:a}
\textit{(a) $\alpha^k:M_0\to H_k^0(M)$ is bijective if $r\geq1$.}

\phantomsection\label{fac:III:3:no:62:prop:1:item:b}
\textit{(b) $H_k^q(M)=0$ for $0<q<r$.}

(For $i=0$, the hypothesis means that $t_0^k.m=0$ implies $m=0$.)

This proposition is a special case of a result due to de Rham [11]
(de Rham's result remains valid even when the
$m\langle i_0\cdots i_q\rangle$ are not assumed homogeneous). See also
[6], Chapter VIII, \S4, where a special case sufficient for our
applications is treated.

We apply \hyperref[fac:III:3:no:62:prop:1]{Proposition~1} to the graded
$S$-module $S(n)$.

\phantomsection\label{fac:III:3:no:62:prop:2}
\textsc{Proposition 2.} \textit{Let $k$ be an integer $\geq0$, and let
$n$ be any integer. Then:}

\phantomsection\label{fac:III:3:no:62:prop:2:item:a}
\textit{(a) $\alpha^k:S_n\to H_k^0(S(n))$ is bijective if $r\geq1$.}

\phantomsection\label{fac:III:3:no:62:prop:2:item:b}
\textit{(b) $H_k^q(S(n))=0$ for $0<q<r$.}

\phantomsection\label{fac:III:3:no:62:prop:2:item:c}
\textit{(c) $H_k^r(S(n))$ has as a basis over $K$ the cohomology
classes of the monomials $t_0^{\alpha_0}\cdots t_r^{\alpha_r}$ such
that $0\leq\alpha_i<k$ and
$\sum_{i=0}^{i=r}\alpha_i=k(r+1)+n$.}

It is clear that the graded $S$-module $S(n)$ satisfies the hypotheses
of \hyperref[fac:III:3:no:62:prop:1]{Proposition~1}; this proves
\hyperref[fac:III:3:no:62:prop:2:item:a]{(a)}
\hyperref[fac:III:3:no:62:prop:2:item:b]{and (b)}. On the other hand, for
every graded $S$-module $M$,
\phantomsection\label{fac:III:3:no:62:formula:top-cohomology-quotient}
\[
H_k^r(M)=
M_{k(r+1)}/(t_0^k,\ldots,t_r^k)M_{kr}.
\]
Now the monomials
\[
t_0^{\alpha_0}\cdots t_r^{\alpha_r},\qquad
\alpha_i\geq0,\qquad
\sum_{i=0}^{i=r}\alpha_i=k(r+1)+n,
\]
form a basis of $S(n)_{k(r+1)}$, and those among them for which at
least one $\alpha_i$ is $\geq k$ form a basis of
$(t_0^k,\ldots,t_r^k)S(n)_{kr}$. This proves
\hyperref[fac:III:3:no:62:prop:2:item:c]{(c)}.

It is convenient to write the exponents $\alpha_i$ in the form
$\alpha_i=k-\beta_i$. The conditions in
\hyperref[fac:III:3:no:62:prop:2:item:c]{(c)} then become
\phantomsection\label{fac:III:3:no:62:formula:beta-conditions}
\[
0<\beta_i\leq k
\qquad\text{and}\qquad
\sum_{i=0}^{i=r}\beta_i=-n.
\]
The second condition, together with $\beta_i>0$, implies
$\beta_i\leq -n-r$. Hence, if $k\geq -n-r$, the condition
$\beta_i\leq k$ follows from the preceding two. Thus:

\phantomsection\label{fac:III:3:no:62:cor:1}
\textsc{Corollary 1.} \textit{For $k\geq-n-r$, $H_k^r(S(n))$ has as
a basis the cohomology classes of the monomials}
\[
\frac{(t_0\cdots t_r)^k}
{t_0^{\beta_0}\cdots t_r^{\beta_r}},
\qquad \beta_i>0,\qquad
\sum_{i=0}^{i=r}\beta_i=-n.
\]

We also have:

\phantomsection\label{fac:III:3:no:62:cor:2}
\textsc{Corollary 2.} \textit{If $h\geq k\geq-n-r$, the homomorphism}
\[
\rho_k^h:H_k^q(S(n))\longrightarrow H_h^q(S(n))
\]
\textit{is bijective for every $q\geq0$.}

For $q\neq r$, this follows from
\hyperref[fac:III:3:no:62:prop:2:item:a]{assertions (a)}
\hyperref[fac:III:3:no:62:prop:2:item:b]{and (b)} of
\hyperref[fac:III:3:no:62:prop:2]{Proposition~2}. For $q=r$, it follows
from \hyperref[fac:III:3:no:62:cor:1]{Corollary~1}, since $\rho_k^h$
sends
\[
\frac{(t_0\cdots t_r)^k}{t_0^{\beta_0}\cdots t_r^{\beta_r}}
\quad\text{to}\quad
\frac{(t_0\cdots t_r)^h}{t_0^{\beta_0}\cdots t_r^{\beta_r}}.
\]

\phantomsection\label{fac:III:3:no:62:cor:3}
\textsc{Corollary 3.} \textit{The homomorphism
$\alpha:S_n\to H^0(S(n))$ is bijective if $r\geq1$ or if $n\geq0$.
We have $H^q(S(n))=0$ for $0<q<r$, and $H^r(S(n))$ is a vector space
of dimension $\binom{-n-1}{r}$ over $K$.}

The assertion concerning $\alpha$ follows from
\hyperref[fac:III:3:no:62:prop:2:item:a]{Proposition~2(a)} when
$r\geq1$; it is immediate when $r=0$ and $n\geq0$. The rest of the
corollary follows at once from
\hyperref[fac:III:3:no:62:cor:1]{Corollaries~1} and
\hyperref[fac:III:3:no:62:cor:2]{2}, with the convention that the
binomial coefficient $\binom{a}{r}$ is zero when $a<r$.
